\documentclass[10pt,a4paper]{amsart}
\usepackage[margin=19mm]{geometry}
\usepackage{amsmath,amssymb,mathtools}
\usepackage[scr=rsfs,cal=euler]{mathalfa}
\usepackage{xfrac}
\usepackage[unicode,hidelinks,bookmarks=false]{hyperref}
\newcommand{\ie}{i.e.\ }
\newcommand{\cf}{cf.\ }
\newcommand{\resp}{respectively\ }
\newcommand{\Subsection}{\subsection}
\providecommand{\nobreakdash}{\nobreak\hbox{-}}
\setlength{\emergencystretch}{2em}
\multlinegap0pt
\usepackage{bm}
\usepackage{bbm}
\usepackage{dsfont}
\usepackage{xspace}
\usepackage{cancel}
\usepackage{mathtools}
\usepackage{amsthm}
\makeatletter
\@ifundefined{thm}{\newtheorem{thm}{Thm}}{}
\@ifundefined{lem}{\newtheorem{lem}[thm]{Lemma}}{}
\@ifundefined{prop}{\newtheorem{prop}[thm]{Proposition}}{}
\makeatother
\title[Simplicial arrangements with few double points]{Simplicial arrangements with few double points}
\author{Dmitri Panov, Guillaume Tahar}
\date{Source: arXiv:2409.01892v2 (CC BY 4.0)}
\begin{document}
\maketitle

\noindent\emph{Source and licence.} Simplicial arrangements with few double points. Version arXiv:2409.01892v2 (CC BY 4.0), distributed under the Creative Commons Attribution 4.0 International licence, as recorded in the arXiv licence metadata for this exact version and shown on the abstract page https://arxiv.org/abs/2409.01892v2. Full rights notice: licenses/arxiv-2409.01892-CC-BY-4.0.txt.\par\smallskip

\noindent\textit{Editorial scope.} Counted unit: the Lem whose source label is \texttt{lem:rigidity1} in the article (source0.tex lines 580--582, source label \texttt{lem:rigidity1}), delivered together with the article's own complete original proof of it (source0.tex lines 584--610, one \texttt{proof} environment of 27 lines). No mathematics is written, repaired or re-derived: statement and proof are the article's own text line for line. Required context retained uncounted: lem:adjaDP (lines 244--246); prop:configuration (lines 523--535); lem:rigidity (lines 564--566). Every definition, display and label of those lines is retained unchanged. Omitted from this package and not redistributed: 1--243, 247--522, 536--563, 567--579, 583--583, 611--760. Nothing inside a retained excerpt is omitted or altered: every retained body is delivered exactly as the article's lines are written, so no omission receipt of any kind accompanies this package. Citations: the retained excerpts carry no citation command at all, so no bibliography entry of the article is transcribed and no reference is redistributed. Reference adaptation: none was needed and none was made; every reference inside the delivered unit resolves inside it, so no unresolved reference or citation is produced. Printed slips kept literal: no printed word, symbol, hypothesis or number of the counted statement or of its proof was altered. Layout and class: the article's own class is \texttt{amsart} with packages including amsthm, amsmath, enumerate, amsfonts, amssymb, fullpage, amscd, stmaryrd, graphicx, tikz-cd, array, dsfont, inputenc, babel; the wrapper is the lane's pinned \texttt{amsart} editorial wrapper at 10pt on A4, which restates the theorem-like environments the retained text needs and adds the article's own macro definitions that the retained text needs (none), with extra wrapper packages (bm, bbm, dsfont, xspace, cancel, mathtools). The delivered numbering is the unit's own. Assets: no retained span contains a figure environment, a graphics-inclusion command, a PDF-page inclusion, a table or a TikZ picture; no image file is copied into this package, the package declares no asset dependency and no image file is redistributed with it. Licence: version arXiv:2409.01892v2 of this article is distributed under the Creative Commons Attribution 4.0 International licence, as recorded in the arXiv licence metadata for this exact version and shown on the abstract page https://arxiv.org/abs/2409.01892v2. A full rights notice is delivered at licenses/arxiv-2409.01892-CC-BY-4.0.txt. Characters: every retained excerpt is pure ASCII, so the delivered engine, XeLaTeX with its default Latin Modern text font, reports no missing character.
\begin{lem}\label{lem:adjaDP}
A simplicial arrangement contains two adjacent double points if and only if it belongs to family $\mathcal{R}(0)$.
\end{lem}

\begin{prop}\label{prop:configuration}
Consider the simplicial arrangement $\mathcal{A}_{r}$ formed by the sides and symmetric axes of a regular $m$-gon $\mathcal{P}$ for $m \geq 3$ ($\mathcal{A}_{r}$ belongs to the combinatorial class $\mathcal{A}(2m,1)$). If $\Lambda$ is a line such that:
\begin{itemize}
    \item $\Lambda \notin \mathcal{A}_{r}$;
    \item $\Lambda$ does not contain the centre $C$ of the regular $m$-gon $\mathcal{P}$;
    \item $\Lambda$ is not the line at infinity (for the Euclidean structure where $\mathcal{P}$ is a regular polygon);
\end{itemize}
then $\Lambda$ does not contain a configuration of five consecutive intersection points $N_{1},N_{2},N_{3},N_{4},N_{5}$ with $\mathcal{A}_{r}$ such that:
\begin{itemize}
    \item $N_{1},N_{3},N_{5}$ are triple points of $\mathcal{A}_{r}$;
    \item $N_{2},N_{4}$ are interior points of edges of $\mathcal{A}_{r}$.
\end{itemize}
\end{prop}

\begin{lem}\label{lem:rigidity}
Consider a simplicial arrangement $\mathcal{A}$ of $n \geq 15$ lines in $\mathbb{RP}^{2}$ that differs by at most $k < \frac{1}{7}n + \frac{1}{7}$ lines from a regular simplicial arrangement $\mathcal{A}_{r}$ formed by the sides and symmetry axes of a regular $m$-gon $\mathcal{P}$ ($\mathcal{A}_{r}$ belongs to the combinatorial class $\mathcal{A}(2m,1)$). Then no line $\Lambda$ of $\mathcal{A} \setminus \mathcal{A}_{r}$ contains the centre $C$ of the regular polygon $\mathcal{P}$. Besides, $\mathcal{A}$ cannot be a near-pencil.
\end{lem}

\begin{lem}\label{lem:rigidity1}
Consider a simplicial arrangement $\mathcal{A}$ of $n \geq 15$ lines in $\mathbb{RP}^{2}$ that differs by at most $k < \frac{1}{13}n - \frac{14}{13}$ lines from a regular simplicial arrangement $\mathcal{A}_{r}$ formed by the sides and symmetric axes of a regular $m$-gon $\mathcal{P}$ ($\mathcal{A}_{r}$ belongs to the combinatorial class $\mathcal{A}(2m,1)$). Then $\mathcal{A} \setminus \mathcal{A}_{r}$ contains at most one line which is the line at infinity (for an Euclidean structure where regular polygon $\mathcal{P}$ is inscribed in a circle).
\end{lem}

\begin{proof}
Consider a line $\Lambda \in \mathcal{A} \setminus \mathcal{A}_{r}$. Conditions of Lemma~\ref{lem:rigidity1} are more restrictive than that of Lemma~\ref{lem:rigidity} so we can assume that $\Lambda$ does not contain the center $C$. We also assume by contradiction that $\Lambda$ is not the line at infinity for an Euclidean structure where regular polygon $\mathcal{P}$ is inscribed in a circle.
\par
Let $M_{1},\dots,M_{t}$ be the cyclical ordered subset of $\Lambda$ consisting of:
\begin{itemize}
    \item[a)] the intersection points of $\Lambda$ with the lines of the symmetric difference $\mathcal{A} \Delta \mathcal{A}_{r}$;
    \item[b)] the double points of arrangement $\mathcal{A}_{r}$ possibly contained in $\Lambda$.
\end{itemize}
%These points are ordered according to the cyclic order of $\Lambda$. 
Let's show first that $t \leq k+1$.  Indeed, line $\Lambda$ has at most $k-1$ intersection points with the lines of $\mathcal{A} \Delta \mathcal{A}_{r}$, so there are at most $k-1$ points of type a) among $M_i$'s. And at most $2$ of $M_i$'s are of type b), since  $\mathcal A_r$ has exactly $2m$ double points, namely the midpoints of edges of the convex polygon $\mathcal{P}$.
%Besides, there is exactly one double point in the interior of each exterior edge of convex polygon $\mathcal{P}$ (and no double point elsewhere). 
%Therefore, $\Lambda$ cannot contain more than two of them. It follows that $t \leq k+1$. 

Next, we claim that there are at most $4t$ lines of ${\mathcal A}\setminus \Lambda$ that are incident to $M_1,\ldots, M_t$. Indeed, since $C\notin \Lambda$, each $M_i$ is at most a triple point of $\mathcal A_r$,  and $|\mathcal{A} \Delta \mathcal{A}_{r}\setminus \Lambda|=k-1$ . Taking the sum of contributions, and using   $t \leq k+1$ we get $3t+k-1\le 4t$.
 
%(indeed, a triple point of $\mathcal{A}_{r}$ can still be one of the $M_{1},\dots,M_{t}$ if it is an intersection point with one line of $\mathcal{A} \Delta \mathcal{A}_{r}$).
\par
Our goal will now be to prove that each open interval $]M_{i},M_{i+1}[$ of $\Lambda$ intersects at most $9$ lines of $\mathcal{A}\setminus \Lambda$. If the set $\lbrace{ M_{1},\dots,M_{t} \rbrace}$ is empty, we will prove similarly that the whole projective line intersects at most $9$ lines of $\mathcal{A}\setminus \Lambda$. Then, counting the lines of $\mathcal{A}\setminus\Lambda$ intersected by $\Lambda$ inside and outside of $M_{1},\dots,M_{t}$, we obtain that $n-1 \leq 4t+9t \leq 13k+13$ (since $t \leq k+1$). This contradicts the inequality $k < \frac{1}{13}n - \frac{14}{13}$ and $\Lambda$ is therefore the line at infinity.\newline

It remains to prove that in each open interval $]M_{i},M_{i+1}[$ (or the whole line if $\lbrace{ M_{1},\dots,M_{t} \rbrace}$ is empty), $\Lambda$ intersects at most $9$ lines of $\mathcal{A} \setminus \lbrace{ \Lambda \rbrace}$. We assume by contradiction that there is such an open interval (that can be the whole line) $]M_{i},M_{i+1}[$ in $\Lambda$ that intersects at least $10$ lines of $\mathcal{A}$. Since the intersections of $\Lambda$ with the lines of $\mathcal{A} \Delta \mathcal{A}_{r}$ are contained in $M_{1},\dots,M_{n}$, all these lines of $\mathcal{A}$ that $\Lambda$ intersects in $]M_{i},M_{i+1}[$ actually are lines of  $\mathcal{A}_{r}$. It also follows from the definition of $M_{1},\dots,M_{t}$ that every line of $\mathcal{A}_{r}$ intersecting $\Lambda$ outside of $M_{1},\dots,M_{t}$ actually belongs to $\mathcal{A}$.
\par
The only vertices of $\mathcal{A}_{r}$ that line $\Lambda$ can meet outside of $M_{1},\dots,M_{t}$ are triple points of $\mathcal{A}_{r}$. Therefore, the vertices of $\mathcal{A}$ contained in the open segment $]M_{i},M_{i+1}[$ are either quadruple points of $\mathcal{A}$ (if they are vertices of $\mathcal{A}_{r}$) or double points of $\mathcal{A}$ (otherwise).
\par
Arrangement $\mathcal{A}$ does not belong to family $\mathcal{R}(0)$ (see the last claim in Lemma~\ref{lem:rigidity}). Hence, there cannot be two consecutive double points in $]M_{i},M_{i+1}[$ (see Lemma~\ref{lem:adjaDP}). We prove now that there cannot be two consecutive quadruple points $A,B$ of $\mathcal{A}$ in $]M_{i},M_{i+1}[$. Indeed, such segment $[A,B] \subset ]M_{i},M_{i+1}[$ is drawn between two vertices of $\mathcal{A}_{r}$ so unless it is contained in a line of $\mathcal{A}_{r}$ (and $\Lambda \notin \mathcal{A}_{r}$), $S$ intersects the interior of an edge of $\mathcal{A}_{r}$. This is a contradiction with the statement that $A,B$ are consecutive intersection points with $\mathcal{A}_{r}$ in $]M_{i},M_{i+1}[$. So vertices of segment $]M_{i},M_{i+1}[$ alternate between double and quadruple points. Since $]M_{i},M_{i+1}[$ intersects at least $10$ lines of $\mathcal{A}_{r}$, we conclude that the open segment $]M_{i},M_{i+1}[$ contains a sequence of five consecutive vertices $N_{1},N_{2},N_{3},N_{4},N_{5}$ where $N_{1},N_{3},N_{5}$ are quadruple points while $N_{2},N_{4}$ are double points.
\par
Since $N_{1},N_{3},N_{5}$ are triple points in $\mathcal{A}_{r}$ while $N_{2},N_{4}$ are interior points of edges in $\mathcal{A}_{r}$, Proposition~\ref{prop:configuration} proves that there is no line containing this configuration of five points. This final contradiction ends the proof.
\end{proof}

\end{document}
